\section{Methodology} \label{sec:method}
% ============================================================================
\subsection{Initial Burst Selection} \label{sec:sel}
We select single-pulse bursts from the Fermi/GBM catalog
\citep{Meegan2009,Gruber2014} using the shape-based ``horizontal-line'' algorithm
of \citet{Busby2024}, which scores how well a single FRED-like (fast-rise,
exponential-decay) pulse describes the background-subtracted light curve. We retain bursts above the \citeauthor{Busby2024}
single-pulse threshold (score $\gtrsim0.9983$) and apply a fluence cut of
$>$$10^{-5}$~erg~cm$^{-2}$ (10--1000~keV, GBM catalog) to ensure adequate
time-resolved signal-to-noise, since a robust spectral shape requires
sufficient counts per bin (Section~\ref{sec:bin}). No $T_{90}$ restriction is
imposed; the sample therefore contains both short- and long-duration bursts whose
emission is dominated by a single pulse. We emphasize that this selection is on
pulse \emph{shape}, not brightness: the sample fluence spans a factor of
$\sim$$250$. The final sample comprises \textbf{$106$} bursts
(Table~\ref{tab:sample}).

\subsection{Detector, Source, and Background Selections} \label{sec:bkg}
For each burst the detector--source angle $\theta$ of the 14 GBM detectors is
computed from the spacecraft attitude (the POSHIST quaternions interpolated to
the trigger time). We use the time-tagged event (TTE) data of the NaI detectors
with $\theta\le50^\circ$; a detector at $50^\circ<\theta\le60^\circ$ is also kept
if it appears in the trigger's on-board (BCAT) detector mask---it triggered, so
the source lay in its field of view---and if no NaI then qualifies the single
closest BCAT NaI is retained. Up to the three lowest-angle NaI are used, each
paired with the BGO on its side of the spacecraft (b0 for n0--n5, b1 for
n6--nb). For the $11$ bursts with usable LAT Low-Energy (LLE) data (a subset of
the $18$ LAT-detected bursts) we add the LLE data above 30~MeV
\citep{Atwood2009}. Each detector's response matrix is used; for multi-epoch
\texttt{rsp2} responses the matrix covering the source interval is taken. We use 8--900~keV for NaI, masking the 33--40~keV region around the iodine
K-edge---a sharp jump in the crystal's absorption efficiency at 33.17~keV
that is difficult to calibrate \citep{Ronchi2020}---with 0.3--40~MeV for BGO
and 30--100~MeV for LLE. Inter-detector calibration is handled by free
effective-area constants in the range 0.8--1.2, frozen to unity for the
reference NaI \citep[as in][]{Ronchi2020}. The source interval is set to bracket the burst emission
(Section~\ref{sec:bin}). Each selected detector receives its own pre- and
post-burst background intervals ($\sim$50--150~s each), fit with a time
polynomial of order $\le$3---the order chosen by a likelihood-ratio test---and
interpolated beneath the burst: the detector background varies smoothly with the
spacecraft's orbital position, so windows on either side of the burst constrain
its level during the burst itself. An accurate background is the prerequisite
for recovering any sub-dominant spectral component, so every window is reviewed
and approved---by a human at an interactive light-curve display, or by a
vision-capable agent---with each approval recorded together with its origin, so
the provenance of every selection is auditable.\footnote{The numbers in this
paper use a provisional, automatically-selected background catalogue pending
completion of this approval pass and the authoritative re-fit; the qualitative
conclusions are not expected to be sensitive to the refinement.}

\subsection{Light-curve Binning} \label{sec:bin}
Within the approved emission interval of each burst, we define the time bins
from the reference NaI (the lowest-angle approved NaI) using the Bayesian Blocks
algorithm \citep{Scargle2013} applied to its background-subtracted 8--900~keV
count rate (the \texttt{3ML} implementation, with the polynomial background
included and a false-alarm prior $p_0=0.01$). Bayesian Blocks partitions a light
curve at the points where the count rate changes significantly, producing bins
that are long during quiet stretches and short across sharp variations, so that
each bin samples an approximately constant spectral state. We then impose a
per-block significance floor: each block's net significance is evaluated with
the modeled-background (Gaussian) Li \& Ma statistic, and blocks below $5\sigma$
are dropped at the leading and trailing edges and merged into a neighbour in the
interior until every surviving block clears $5\sigma$, so that each time bin is
simultaneously a distinct count-rate state and a statistically adequate
spectrum. The resulting block edges are shared across all of the burst's
detectors. A dedicated fine-grid pass (1~ms cells with the Poisson block
fitness, run on each burst's pulse core) shows that only $10\%$ ($11/106$) of
the sample contains significant sub-$0.1$~s structure---most prominently
GRB~130427A, with credible blocks down to 2~ms---so intra-bin spectral evolution
is a quantified, per-burst caveat for that fast minority rather than a
sample-wide one. The provisional run yields \textbf{$1057$} time-resolved
spectra, to be regenerated by the authoritative re-block on the approved
backgrounds.

For each block, the spectral data are assembled in \texttt{3ML}: per detector, a
source count spectrum is extracted over the block interval, the background count
spectrum is the polynomial model (Section~\ref{sec:bkg}) integrated over the same
interval with its Gaussian uncertainty, and the corresponding detector response
matrix folds the model into count space. The detectors of a burst are then fit
jointly over their respective energy ranges (Section~\ref{sec:fit}); each block
is one such joint fit.

\subsection{Block Definition and Grading} \label{sec:grade}
We grade each burst by the quality of its time-resolved coverage, following a
three-tier scheme:
\begin{enumerate}
\item \emph{Gold}: $40$ bursts with $\ge10$ well-constrained
  bins, suitable for per-burst parameter relations.
\item \emph{Silver}: $44$ bursts with $5$--$9$ bins, contributing to
  population statistics with caveats.
\item \emph{Bronze}: $22$ bursts with $\le4$ bins, too few for per-burst
  trends; included only in the pooled distributions.
\end{enumerate}
The grade is carried as a column in Table~\ref{tab:perburst} and is used in
Section~\ref{sec:discussion} to show that our main results hold on the cleanest
(Gold) subsample.

\subsection{Spectral Models} \label{sec:models}
We fit each spectrum with six photon models spanning the two interpretations of
the spectral curvature while remaining empirical and comparable with the GBM
catalogs \citep[e.g.,][]{Kaneko2006,Gruber2014}. Empirical functions are flexible
and model-agnostic; although they do not map uniquely onto a mechanism, the
low-energy index $\alpha$ still discriminates between scenarios at high
significance \citep[e.g.,][]{Acuner2020}.

The Band function \citep{Band1993} is a smoothly-joined broken power law; below
the break $E_b=(\alpha-\beta)\Ep/(2+\alpha)$,
\begin{equation}
N(E)=A\,(E/E_0)^{\alpha}\,\exp\!\left[-(2+\alpha)E/\Ep\right],
\end{equation}
joining $N(E)\propto E^{\beta}$ above $E_b$, with $E_0=100$~keV. The cutoff power
law (CPL), $N(E)=A\,(E/E_0)^{\Gamma}\exp(-E/E_c)$ with $\Ep=(2+\Gamma)E_c$, is the
$\beta\!\to\!-\infty$ limit. The smoothly broken power law (SBPL) joins two
power-law segments with photon indices $\alpha$ and $\beta$ at a break energy
$E_{\rm b}$, with the sharpness of the join an adjustable parameter; we use
the GBM-catalog parameterization \citep{Kaneko2006} as implemented in
\textsf{astromodels}, with the break scale held fixed.

To represent the synchrotron interpretation we use the double smoothly broken
power law \citep[2SBPL;][whose Equations~1--4 give the functional
form]{Ravasio2018}: three power-law segments with photon indices
$\alpha_1,\alpha_2,\beta$ joined at two break energies---a lower break $x_b$
and an upper break $x_p$ (the latter at the $\nu F_\nu$ peak)---each join
having a sharpness parameter ($n_1$, $n_2$). We fit the
\textsf{astromodels} implementation, whose free parameters are the
normalization, the three indices, and the two break energies. Under the
optically-thin synchrotron interpretation the two breaks are
physical---$x_b$ is the cooling frequency \nuc\ and $x_p=\num$ the injection
frequency---and the segment between them carries the synchrotron slopes
$\alpha_1\simeq-2/3$ (slow cooling) steepening toward $-3/2$ (fast cooling).
These characteristic values follow from the radiation physics itself: the
low-energy tail of even a single electron's synchrotron spectrum cannot be
harder than $N(E)\propto E^{-2/3}$, and a fully cooled electron population
softens it to $E^{-3/2}$, so measured indices above these ``lines of death''
challenge a pure synchrotron origin
\citep[][]{Preece1998,Sari1998,Ravasio2019}. We hold the break sharpnesses fixed
($n_1=n_2=2$ in this provisional run; the verified re-fit will adopt the
$n_1=5.38$, $n_2=2.69$ calibration of \citealt{Ravasio2018}, who found that
leaving $n_1$ free does not converge reliably---the same behavior we observe).
The full fit configuration (parameter bounds, seeds, and restart grids) is
listed in Appendix~\ref{app:fitconfig}.

To represent the photospheric interpretation we add a blackbody to the two
simplest non-thermal continua, following the multi-component decomposition of
\citet{Guiriec2011,Guiriec2015}, in which a sub-dominant quasi-thermal component
accompanies the non-thermal emission and the two are fit \emph{simultaneously}.
The blackbody photon spectrum is $N(E)\propto E^{2}/[\exp(E/\kT)-1]$, and we form
the composites Band$+$BB and CPL$+$BB; we quantify the thermal contribution by the
flux ratio $F_{\rm BB}/F_{\rm tot}$ \citep{Guiriec2015}. A BB$+$SBPL composite
reproduces the curvature of a 2SBPL near the peak, so the $+$BB models act as a
practical thermal proxy for the two-break model.

\subsection{Spectral Fitting and Model Selection} \label{sec:fit}
All fits are performed in the Multi-Mission Maximum Likelihood framework
\citep[3ML;][]{Vianello2015} with the \texttt{pgstat} likelihood, the
appropriate statistic when the source counts are Poisson-distributed and the
background level carries Gaussian uncertainties from the polynomial fit. Model
comparison uses the Akaike and Bayesian information criteria
\citep[AIC, BIC;][]{Akaike1974,Schwarz1978} computed from the standard
$-2\ln L$. An information criterion is the fit's $-2\ln L$ penalized by the
number of free parameters, so a lower value indicates a better description of
the data after accounting for model complexity; a gap of $\daic\ge10$ between
two models is conventionally decisive. Which statistic is \emph{valid},
however, depends on the relationship between the models compared. For \emph{nested} pairs---where the simpler model is a special case of the
more complex one, as for Band$+$BB/Band, CPL$+$BB/CPL, and 2SBPL/SBPL---a
blackbody or second break is deemed significant at $\daic\ge10$ relative to the
parent continuum---the AIC analog of the $\Delta\mathrm{DIC}\ge10$ threshold of
\citet{Li2021} and the Bayes-factor threshold ($2\ln K>10$) of
\citet{Burgess2019}, and equivalent here to a likelihood-ratio statistic
\citep{Wilks1938} of $\ge$$14$, beyond the nominal 99\% (2~d.o.f.) point; we
prefer the information-criterion form because the $\chi^2$ calibration of the
classical LRT is only approximate when a component normalization sits at a
parameter boundary. Every blackbody and second break that passes $\daic\ge10$
also passes the more conservative $\Delta\mathrm{BIC}$ criterion. For
\emph{non-nested} comparisons---in particular the central thermal
($+$BB) versus two-break (2SBPL) question, where neither model contains the
other as a limiting case---the LRT is undefined, and we rely
on \daic\ alone, with $\Delta\mathrm{BIC}$ as a cross-check. A physical-validity
gate prevents a fit from \emph{winning} the selection if any key parameter is
railed against a bound or, for the 2SBPL, if the breaks are mis-ordered
($x_b\ge x_p$); in the $4\%$ ($43/1057$) of bins where no model passes the gate
we assign no preferred model. The blackbody possesses a railing local minimum,
so we multi-start the $+$BB models whose blackbody is railed or insignificant
from a grid of temperature seeds and retain the lowest $-2\ln L$. The 2SBPL has
no analogous restart in the present (provisional) run, so two-break counts are
conservative lower limits (Section~\ref{sec:curv}). We quantify each parameter
correlation by the Spearman rank coefficient $\rho$ and a least-squares slope in
log space. Energy and photon fluxes are obtained by integrating the best-fit
Band model over 10--1000~keV, with uncertainties propagated by Monte Carlo over
the fitted parameter uncertainties (200 draws; 16th--84th percentiles). For the
principal $\num$--$\nuc$ relation, whose two variables carry comparable
uncertainties, we maximize the \citet{DAgostini2005} errors-in-both-variables
likelihood,
\begin{equation}
-2\ln\mathcal{L}=\sum_i\left[\ln 2\pi\sigma_i^2
+\frac{(y_i-m x_i-c)^2}{\sigma_i^2}\right],\quad
\sigma_i^2=\ssc^2+\sigma_{y,i}^2+m^2\sigma_{x,i}^2,
\label{eq:dagostini}
\end{equation}
with the measured break uncertainties (in dex) on both axes and the intrinsic
scatter \ssc---the spread of the relation beyond what the measurement errors
alone explain---a free parameter.

% ============================================================================
